% chapters/glossary.tex — 용어집 (알파벳순)

\chapter*{용어집}
\addcontentsline{toc}{chapter}{용어집}
\markboth{용어집}{용어집}

이 용어집은 본 책에서 사용하는 주요 용어를 알파벳순으로 정리한 것이다. 각 항목은 짧은 정의와 함께 해당 장 번호를 표시하여 독자가 쉽게 참조할 수 있도록 구성했다. 처음 책을 읽을 때는 가볍게 훑어보고, 본문을 읽다가 모르는 용어가 나오면 다시 돌아와 참조하기 바란다.

\begin{description}[leftmargin=2.5em, style=nextline, font=\bfseries\sffamily\color{inkblue}]

\item[Activation (활성값)]
신경망의 각 층을 통과한 후 출력되는 값. 레이어의 출력 텐서를 의미하며, 다음 층의 입력이 된다. \textit{참고: 4장}

\item[AdamW]
Adam 옵티마이저에 weight decay를 결합한 최적화 알고리즘. 현대 LLM 학습의 사실상 표준. \textit{참고: 7장}

\item[Attention (어텐션)]
시퀀스의 각 위치가 다른 모든 위치에 얼마나 주의를 기울일지 학습하는 메커니즘. 트랜스포머의 핵심. \textit{참고: 5장}

\item[Backpropagation]
신경망의 손실을 각 파라미터로 미분하여 그래디언트를 계산하는 알고리즘. 연쇄 법칙(chain rule)에 기반. \textit{참고: 7장}

\item[BPE (Byte Pair Encoding)]
바이트 단위로 시작하여 가장 자주 등장하는 쌍을 병합해 나가는 토크나이저 알고리즘. GPT-2, GPT-3 등이 사용. \textit{참고: 3장}

\item[Causal Mask (인과 마스크)]
자기회귀 모델에서 미래 토큰을 보지 못하도록 가리는 하삼각 마스크. \textit{참고: 5장}

\item[Context Length (컨텍스트 길이)]
모델이 한 번에 처리할 수 있는 최대 토큰 수. GPT-2는 1024, GPT-4는 128K. \textit{참고: 7장}

\item[Cross-Entropy Loss (교차 엔트로피 손실)]
분류 문제에서 예측 확률 분포와 정답 분포 사이의 차이를 측정하는 손실 함수. 언어 모델링의 표준 손실. \textit{참고: 7장}

\item[Decoder (디코더)]
트랜스포머에서 시퀀스를 생성하는 부분. GPT는 디코더만 사용하는 구조. \textit{참고: 6장}

\item[Dropout]
학습 시 일부 뉴런을 무작위로 0으로 만들어 과적합을 방지하는 정규화 기법. \textit{참고: 4장}

\item[Embedding (임베딩)]
이산적인 토큰 ID를 연속적인 밀집 벡터로 변환하는 과정. 또는 그 벡터 자체. \textit{참고: 4장}

\item[Feed-Forward Network (FFN, 순방향 신경망)]
트랜스포머 블록에서 어텐션 다음에 적용되는 위치별 완전연결 신경망. 보통 2개의 선형 레이어 + 활성화 함수. \textit{참고: 6장}

\item[GELU (Gaussian Error Linear Unit)]
ReLU의 부드러운 버전. 가우시안 함수를 사용하여 0 주변에서 부드럽게 전환. GPT 모델이 사용. \textit{참고: 6장}

\item[Gradient (그래디언트)]
손실 함수의 각 파라미터에 대한 미분값. 최적화 방향을 가리킴. \textit{참고: 7장}

\item[Greedy Decoding (탐욕 디코딩)]
각 스텝에서 가장 확률이 높은 토큰만 선택하는 생성 전략. 빠르지만 단조로운 출력. \textit{참고: 8장}

\item[Head (헤드)]
멀티헤드 어텐션에서 독립적으로 어텐션을 수행하는 단위. \textit{참고: 5장}

\item[Layer Normalization (레이어 정규화)]
텐서의 마지막 차원을 따라 평균과 분산을 정규화하는 기법. 트랜스포머 학습 안정화에 필수. \textit{참고: 4장, 6장}

\item[Learning Rate (학습률)]
그래디언트의 스텝 크기를 결정하는 하이퍼파라미터. 너무 크면 발산, 너무 작으면 느린 수렴. \textit{참고: 7장}

\item[Logits (로짓)]
softmax를 적용하기 전의 원시 출력 값. 보통 모델의 마지막 선형 레이어 출력. \textit{참고: 7장}

\item[Loss Function (손실 함수)]
모델의 예측과 정답 사이의 차이를 스칼라로 측정하는 함수. 학습의 최적화 대상. \textit{참고: 7장}

\item[Multi-Head Attention (멀티헤드 어텐션)]
여러 개의 어텐션 헤드를 병렬로 수행하여 다양한 패턴을 동시에 학습. \textit{참고: 5장}

\item[Parameter (파라미터)]
모델이 학습하는 가중치. 선형 레이어의 W, b 등. \textit{참고: 7장}

\item[Perplexity (퍼플렉서티, 혼란도)]
언어 모델의 평가 지표. $PPL = \exp(\text{loss})$. 낮을수록 좋은 모델. \textit{참고: 7장}

\item[Positional Encoding (위치 인코딩)]
트랜스포머에 순서 정보를 주입하는 방법. sin/cos 또는 학습 가능한 임베딩. \textit{참고: 4장}

\item[Pre-LN / Post-LN]
레이어 정규화의 위치. Pre-LN은 어텐션/FFN 전에, Post-LN은 후에 적용. Pre-LN이 학습 안정성 우수. \textit{참고: 6장}

\item[Residual Connection (잔차 연결)]
입력을 출력에 더하는 연결. $y = x + f(x)$. 깊은 네트워크 학습에 필수. \textit{참고: 6장}

\item[Scaled Dot-Product Attention]
$Q \cdot K^T / \sqrt{d_k}$ 형태의 어텐션. 트랜스포머의 가장 기본 연산. \textit{참고: 5장}

\item[Self-Attention (자기 어텐션)]
같은 시퀀스 내의 토큰들 사이의 어텐션. 트랜스포머의 핵심. \textit{참고: 5장}

\item[Softmax]
벡터를 확률 분포로 변환하는 함수. $\text{softmax}(x_i) = e^{x_i} / \sum_j e^{x_j}$. \textit{참고: 5장}

\item[Temperature (온도)]
샘플링 시 분포의 날카로움을 조절하는 파라미터. 낮을수록 결정론적, 높을수록 무작위. \textit{참고: 8장}

\item[Token (토큰)]
텍스트를 분할한 최소 단위. 단어, 부분단어, 또는 바이트일 수 있음. \textit{참고: 3장}

\item[Tokenizer (토크나이저)]
텍스트를 토큰 ID 시퀀스로 변환하고, 그 반대 과정도 수행하는 모듈. \textit{참고: 3장}

\item[Top-K Sampling]
다음 토큰 후보 중 확률이 가장 높은 K개만 고려하여 샘플링하는 전략. \textit{참고: 8장}

\item[Top-P (Nucleus) Sampling]
누적 확률이 P가 될 때까지의 토큰만 고려하는 샘플링 전략. 분포 형태에 따라 적응적. \textit{참고: 8장}

\item[Transformer (트랜스포머)]
Vaswani et al. (2017)이 제안한 어텐션 기반 신경망 구조. 현대 LLM의 기반. \textit{참고: 6장}

\item[Unigram Tokenizer]
각 토큰에 확률을 할당하고 EM으로 학습하는 토크나이저. SentencePiece가 사용. \textit{참고: 3장}

\item[Viterbi Algorithm]
동적 계획법으로 최적 분할을 찾는 알고리즘. Unigram 토크나이저의 디코딩에 사용. \textit{참고: 3장}

\item[Vocabulary (어휘)]
토크나이저가 인식하는 모든 토큰의 집합. 어휘 크기는 보통 1만~10만. \textit{참고: 3장}

\item[Weight Tying (가중치 타이)]
입력 임베딩과 출력 레이어의 가중치를 공유하여 파라미터 수를 줄이는 기법. \textit{참고: 7장}

\end{description}
